\documentclass[a4paper,11pt]{article}
\usepackage[ngerman]{babel}
\usepackage[T1]{fontenc}
\usepackage[utf8]{inputenc}
\usepackage[a4paper,margin=2.5cm]{geometry}
\usepackage{microtype}
\usepackage{fancyhdr}
\usepackage{csquotes}
\usepackage{lmodern}
\usepackage{tcolorbox}
\usepackage{enumitem}
\usepackage{parskip}
\usepackage{graphicx}
\usepackage{float}
\usepackage{booktabs}
\usepackage{array}
\usepackage{plantuml}

\pagestyle{fancy}
\geometry{margin=1.5cm}

\begin{document}
	
	\fancyhf{}
	\fancyhead[L]{Jonas Fuchs 1603814}
	
	\section{Aufgabe}
	\subsection*{c)}
	
	Bei der Verwendung von \texttt{notify()} kann es zu Problemen kommen, da immer nur ein wartender Thread aufgeweckt wird. Dabei ist nicht sichergestellt, dass der richtige Thread ausgewählt wird.
	
	Zum Beispiel kann ein Producer aufgeweckt werden, obwohl eigentlich ein Consumer benötigt wird. Der Producer merkt dann, dass der Stack immer noch voll ist, und muss wieder warten. Währenddessen warten andere Threads weiterhin, obwohl sie eigentlich ihre Aufgabe ausführen könnten.
	
	Dadurch kann es zu einem sogenannten Benachrichtigungsstillstand kommen.
	
	Nach Coffman gibt es vier Bedingungen, die gleichzeitig erfüllt sein müssen, damit ein klassischer Deadlock entstehen kann:
	
	\begin{enumerate}
		\item \textbf{Mutual Exclusion (Gegenseitiger Ausschluss):}
		Der Zugriff auf den Stack wird durch einen Lock geschützt. Deshalb kann immer nur ein Thread gleichzeitig auf den kritischen Bereich zugreifen.
		
		\item \textbf{Hold and Wait (Halten und Warten):}
		Ein Thread wartet auf eine bestimmte Bedingung, während andere Threads noch Ressourcen belegen.
		
		\item \textbf{No Preemption (Keine Entziehung):}
		Ein Lock kann einem Thread nicht einfach weggenommen werden. Der Thread muss ihn selbst wieder freigeben.
		
		\item \textbf{Circular Wait (Zyklisches Warten):}
		Mehrere Threads warten gegenseitig auf Ressourcen, die gerade von anderen Threads verwendet werden.
	\end{enumerate}
	
	Wenn alle vier Bedingungen gleichzeitig auftreten, kann ein Deadlock entstehen.
	
	Bei dem vorliegenden Programm muss man allerdings beachten, dass ein Problem mit \texttt{notify()} nicht direkt ein klassischer Deadlock nach Coffman ist. Hier liegt das Problem eher darin, dass möglicherweise der falsche Thread aufgeweckt wird und dadurch andere Threads weiter warten müssen.
	
	
	\section{Aufgabe}
	
	\subsection{}
	
	Ein \textbf{Data Race} entsteht, wenn mehrere Threads gleichzeitig auf dieselbe Speicherstelle zugreifen und mindestens einer davon einen Wert verändert. Wenn die Zugriffe nicht richtig synchronisiert werden, kann es zu falschen oder unerwarteten Ergebnissen kommen.
	
	Eine \textbf{Race Condition} liegt vor, wenn das Ergebnis eines Programms davon abhängt, welcher Thread gerade zuerst ausgeführt wird. Da die Threads unterschiedlich schnell laufen können, ist das Ergebnis nicht immer gleich.
	
	Ein Data Race kann also eine Race Condition verursachen. Die beiden Begriffe bedeuten aber nicht genau dasselbe. Eine Race Condition kann auch auftreten, wenn kein direkter Data Race vorhanden ist, zum Beispiel wenn mehrere Threads in einer ungünstigen Reihenfolge auf gemeinsame Ressourcen zugreifen.
	
	\subsection{}
	
	Das Programm soll die Anzahl der Besucher mithilfe von zehn verschiedenen Threads zählen. Jeder Thread durchläuft eine Schleife mit 1000 Wiederholungen und erhöht dabei die globale Variable \texttt{counter} um eins.
	
	Insgesamt sollte also folgender Wert herauskommen:
	
	\[
	10 \cdot 1000 = 10000
	\]
	
	Das Problem ist, dass alle zehn Threads auf dieselbe Variable \texttt{counter} zugreifen. Dabei wird der aktuelle Wert zuerst in der lokalen Variable \texttt{temp} gespeichert. Danach wird dieser um eins erhöht und wieder in \texttt{counter} geschrieben.
	
	Diese Schritte werden nicht zusammen geschützt. Deshalb können mehrere Threads gleichzeitig auf denselben Wert zugreifen.
	
	Es kann zum Beispiel passieren, dass zwei Threads gleichzeitig den Wert 5 auslesen. Beide erhöhen ihn auf 6 und schreiben diesen Wert anschließend zurück. Eigentlich hätte der Zähler auf 7 steigen müssen, aber durch die gleichzeitigen Zugriffe wird nur eine Erhöhung gezählt.
	
	Der vorhandene \texttt{sleep()}-Aufruf kann das Problem zusätzlich verstärken. Während ein Thread wartet, können andere Threads weiterlaufen und ebenfalls den Zähler verändern.
	
	Dadurch können einzelne Erhöhungen verloren gehen. Am Ende ist die gezählte Besucherzahl deshalb meistens kleiner als die erwarteten 10000.
	
	Da die Threads nicht immer in derselben Reihenfolge ausgeführt werden, kann das Ergebnis bei jedem Programmstart anders sein.
	
	\subsection{}
	
	Um das Problem zu lösen, kann ein Mutex-Lock mit \texttt{threading.Lock()} verwendet werden.
	
	Der kritische Bereich wird dabei mit \texttt{with lock:} geschützt. Dadurch kann immer nur ein Thread gleichzeitig auf den Zähler zugreifen und ihn verändern.
	
	Der Thread muss zuerst den Lock bekommen, bevor er den aktuellen Wert ausliest. Erst nachdem der neue Wert gespeichert wurde, wird der Lock wieder freigegeben.
	
	So wird verhindert, dass zwei Threads gleichzeitig denselben alten Wert auslesen und sich gegenseitig beim Speichern überschreiben.
	
	Der vorhandene \texttt{sleep()}-Aufruf bleibt dabei unverändert. Da er sich innerhalb des geschützten Bereichs befindet, bleibt der Lock auch während dieser Wartezeit belegt. Die anderen Threads müssen deshalb warten, bis der aktuelle Thread fertig ist und den Lock freigibt.
	
	Durch diese Lösung werden die Zugriffe auf den Zähler nacheinander ausgeführt und das Data Race wird verhindert.
	
	Am Ende sollte dadurch immer das erwartete Ergebnis herauskommen:
	
	\[
	10 \cdot 1000 = 10000
	\]
	
	Die gezählte Besucherzahl entspricht somit der tatsächlichen Anzahl der durchgeführten Erhöhungen.
	
\end{document}
